\documentclass[11pt]{article}

\usepackage[T1]{fontenc}
\usepackage[utf8]{inputenc}
\usepackage{lmodern}
\usepackage[a4paper,margin=1in]{geometry}
\usepackage{amsmath,amssymb,amsthm,mathtools}
\usepackage{booktabs}
\usepackage{array}
\usepackage{xcolor}
\usepackage{microtype}
\usepackage[hidelinks]{hyperref}

\setlength{\parindent}{0pt}
\setlength{\parskip}{0.55em}
\setlength{\emergencystretch}{3em}

\newtheorem{lemma}{Lemma}[section]
\newtheorem{theorem}[lemma]{Theorem}
\newtheorem{proposition}[lemma]{Proposition}
\newtheorem{invariant}[lemma]{Invariant}
\theoremstyle{definition}
\newtheorem{definition}[lemma]{Definition}
\theoremstyle{remark}
\newtheorem{remark}[lemma]{Remark}

\newcommand{\Tid}{\mathsf{Tid}}
\newcommand{\Loc}{\mathsf{Loc}}
\newcommand{\Val}{\mathsf{Val}}
\newcommand{\Node}{\mathsf{Node}}
\newcommand{\TS}{\mathsf{TS}}
\newcommand{\topts}{\mathsf{TSTop}}
\newcommand{\interval}[2]{[#1,#2]}
\newcommand{\dom}{\mathop{\mathrm{dom}}}
\newcommand{\fst}{\mathop{\mathrm{fst}}}
\newcommand{\live}{\mathsf{live}}
\newcommand{\pending}{\mathsf{pending}}
\newcommand{\fresh}{\mathsf{fresh}}
\newcommand{\Ready}{\mathsf{Ready}}
\newcommand{\Atomic}{\mathsf{Atomic}}
\newcommand{\Inv}{\mathsf{inv}}
\newcommand{\Lini}{\mathsf{lini}}
\newcommand{\Linr}{\mathsf{linr}}
\newcommand{\GetTop}{\mathsf{getTop}}
\newcommand{\Push}{\mathsf{push}}
\newcommand{\TryRemove}{\mathsf{tryRemove}}
\newcommand{\Succ}{\mathsf{Succ}}
\newcommand{\Fail}{\mathsf{Fail}}
\newcommand{\Empty}{\mathsf{Empty}}
\newcommand{\true}{\mathsf{true}}
\newcommand{\false}{\mathsf{false}}
\newcommand{\arrayedge}{E_{\!a}}
\newcommand{\poolgraph}{E_{\!p}}
\newcommand{\poss}{\Delta}
\newcommand{\conc}{\sigma}
\newcommand{\abst}{\rho}
\newcommand{\tokens}{\pi}
\newcommand{\oplusset}{\mathbin{\oplus}}
\newcommand{\wandsep}{\mathbin{-\!*\!-}}
\newcommand{\entails}{\mathrel{\vDash}}
\newcommand{\tran}{\mathrel{\rightsquigarrow}}
\newcommand{\code}[1]{\texttt{\detokenize{#1}}}
\newcommand{\alltok}[2]{#1\mathrel{\mapsto_{\forall}}#2}
\newcommand{\sometok}[2]{#1\mathrel{\mapsto_{\exists}}#2}
\newcommand{\circinv}[2]{#1\mathrel{\mapsto_{#2}}\circ(\Inv)}
\newcommand{\bulletinv}[2]{#1\mathrel{\mapsto_{#2}}\bullet(\Lini)}
\newcommand{\bulletret}[3]{#1\mathrel{\mapsto_{#2}}\bullet(\Linr,#3)}

\title{A Rely--Guarantee Hoare Proof of the ListPool Layer\\
  \large over SPListArray and Interval Timestamps}
\author{Proof exposition for the mechanized development}
\date{\today}

\begin{document}
\maketitle

\begin{abstract}
This note explains the correctness argument for the ListPool layer whose
underlay is the horizontal composition of an SPListArray and an interval
timestamp object.  The central result is a set-valued rely--guarantee Hoare
proof for the three ListPool methods.  The technically interesting method is
\(\GetTop\): its implementation scans rows non-atomically, retains the maximal
timestamp observed, and uses a second counter read to distinguish a stable
empty pool from interference.  The proof therefore keeps a set of abstract
possibilities, including both an atomic fallback and many interval-snapshot
branches.  We present the representation invariant, the rectangularity
condition that makes independent snapshot choices compositional, the
rely--guarantee relations, the loop invariant, the candidate-commit argument,
and the empty/failure counter argument.  Coq-specific proof engineering is
omitted except where it reflects a genuine mathematical issue.
\end{abstract}

\tableofcontents

\section{Statement of the refinement}

Fix a finite, duplicate-free list of participating threads
\(\Upsilon\subseteq\Tid\).  A node identifier is a pair
\[
  n=(\beta,\ell)\in\Node\triangleq\Tid\times\Loc .
\]
The implementation layer imports
\[
  \mathcal E \;\triangleq\;
  \mathsf{SPListArray}(\Upsilon)\otimes \mathsf{Timestamp}
\]
and exports the ListPool interface
\[
  \mathcal F \;\triangleq\;
  \{\Push(v),\GetTop,\TryRemove(\beta,\ell)\}.
\]
The implementation is
\begin{align*}
\Push_\alpha(v):\quad &
  \ell\leftarrow\mathsf{array.insert}(v);\quad
  t\leftarrow\mathsf{newTS}();\quad
  \mathsf{array.setTS}(\ell,t);\quad \mathsf{return}();
\\
\TryRemove_\alpha(\beta,\ell):\quad &
  b\leftarrow\mathsf{array.tryRemove}(\beta,\ell);\quad
  \mathsf{return}(b).
\end{align*}
For \(\GetTop\), the caller resets its iterator, scans every row
\(\beta\in\Upsilon\), retains the greatest observed timestamp, and accumulates
the counters returned by empty rows.  If a candidate was found it is returned.
Otherwise the implementation reads the aggregate counter once more and
returns \(\Succ(\Empty)\) exactly when this read equals the accumulated count;
it returns \(\Fail\) otherwise.

The verified statement is an implementation simulation
\[
  \mathcal M_{\mathsf{LP}}:\mathcal E\preccurlyeq\mathcal F,
\]
from which contextual linearizability of the layer follows by the generic
simulation-to-linearizability theorem.  The simulation proof is discharged by
one rely--guarantee Hoare triple per exported method.

\section{Concrete and abstract transition systems}

\subsection{Interval timestamps}

Timestamps are either the sentinel \(\topts\), stored in a newly allocated
node, or a closed natural-number interval:
\[
  t ::= \topts \mid \interval{l}{u}, \qquad l\le u.
\]
The strict order is deliberately partial:
\[
\interval{l_1}{u_1}<_{\TS}\interval{l_2}{u_2}
  \quad\Longleftrightarrow\quad u_1<l_2,
  \qquad
\interval{l}{u}<_{\TS}\topts,
  \qquad
\topts\not<_{\TS}t.
\]
Thus overlapping intervals are incomparable.  The timestamp object has a
clock \(\tau\) and a partial map \(q:\Tid\rightharpoonup\mathbb N\) of pending
invocations.  Invocation by \(\alpha\) records \(q(\alpha)=\tau\).  A response
may return \(\interval{l}{u}\) when \(q(\alpha)=l\) and \(l\le u\), after which
the clock becomes at least \(u+1\).  Its invariant is
\[
  q(\alpha)=l \quad\Longrightarrow\quad l\le\tau.
\]
Consequently every interval completed before a later invocation is strictly
below the interval returned by that later invocation.

\subsection{The SPListArray state}

Only the following mathematical view of the array state \(a\) is needed:
\[
  a=(C,V,T,g,O,S,PC).
\]
Here \(C:\Tid\to\mathbb N\) contains monotonically increasing per-row
counters; \(V:\Node\rightharpoonup\Val\) contains every allocated value;
\(T:\Node\rightharpoonup\TS\) stores timestamps; \(g\subseteq\Node\) is the
monotone garbage set; \(O(\beta)\) is the current live order of row \(\beta\);
\(S(\alpha)\) is the scan progress of caller \(\alpha\); and
\(PC(\alpha)\), when defined, is the aggregate counter saved at invocation of
\(\mathsf{getCounter}\).

The concrete timestamp edge points from newer to older:
\[
  \arrayedge(n_1,n_2)
  \quad\Longleftrightarrow\quad
  \exists t_1,t_2.\ T(n_1)=t_1\land T(n_2)=t_2\land t_2<_{\TS}t_1.
\]
Insertion in row \(\alpha\) allocates \((\alpha,\ell)\) with timestamp
\(\topts\), increments \(C(\alpha)\), and prepends \(\ell\) to
\(O(\alpha)\).  Removal adds a node to \(g\) and removes its address from the
live order, but does not erase its value or timestamp.

A row scan saves both \(O(\beta)\) and \(C(\beta)\) at invocation.  Its
response returns either the first still-live address from that saved order or
the saved counter when no such address remains.  The scan-progress state
classifies nodes relative to caller \(\alpha\) as
\(\mathsf{Unvisited}\), \(\mathsf{Visiting}\),
\(\mathsf{Visited}\), or \(\mathsf{Ignored}\).  A node inserted after its row
was snapshotted is \(\mathsf{Ignored}\) for that scan.

\subsection{The abstract ListPool DAG}

An abstract pool state is
\[
  p=(V_p,E_p,S_p,P_p,g_p),
\]
where \(V_p:\Node\rightharpoonup\Val\),
\(E_p\subseteq\Node\times\Node\),
\(S_p:\Tid\rightharpoonup\mathcal P(\Node)\) stores per-thread snapshots,
\(P_p:\Tid\rightharpoonup\Loc\) stores pending pushes, and
\(g_p\subseteq\Node\) is garbage.  Edges point from a newer/higher node to an
older/lower node.  For \(N\subseteq\Node\),
\[
  \mathsf{top}(N,E_p,n)
  \;\triangleq\;
  n\in N\ \land\ \forall n'\in N.\ \neg E_p(n',n).
\]

The relevant interval-sequential transitions are as follows.

\begin{description}
\item[Push invocation.]  For fresh \(n=(\alpha,\ell)\), add
  \(V_p(n)=v\), record \(P_p(\alpha)=\ell\), and add
  \[
    \{(n,n')\mid n'\text{ is live and }n'\notin |P_p|\}.
  \]
  Hence a new push is ordered after every already completed live push, but is
  not ordered against overlapping pending pushes.

\item[Push response.] Remove \(\alpha\) from \(P_p\).

\item[Snapshot invocation.] Set
  \(S_p(\alpha)=\dom(V_p)\), leaving the graph unchanged.

\item[Node response.] If \(S_p(\alpha)=N\), a response
  \(\Succ(v,n)\) is legal when either
  \(n\in g_p\) and \(V_p(n)=v\), or
  \[
    \mathsf{top}(N\setminus g_p,E_p,n)\land V_p(n)=v.
  \]
  The response clears \(S_p(\alpha)\).

\item[Atomic alternatives.] A \(\GetTop\) invocation may instead enter an
  atomic pending state.  From there it may return \(\Fail\) unconditionally,
  or \(\Succ(\Empty)\) when every allocated vertex is garbage.

\item[Removal.] \(\TryRemove(n)\) atomically returns \(\true\) and adds
  \(n\) to \(g_p\) when \(n\) is live; it returns \(\false\) when
  \(n\in g_p\).
\end{description}

Calls by actors outside \(\Upsilon\), and removals with an owner outside
\(\Upsilon\) or an undefined node, have matching abstract error transitions.
This alignment is required for the error branch of each Hoare triple.

\section{Set-valued logical states}

\subsection{Possibilities and operation tokens}

A proof state is
\[
  w=(\conc,\poss),
\]
where \(\conc\) is one concrete underlay state and \(\poss\) is a nonempty set
of abstract possibilities.  A possibility is a pair \((\abst,\tokens)\) of an
abstract ListPool control state and a linearization-token map.  For operation
\(m\), the token of thread \(\alpha\) is one of
\[
  \Inv(m),\qquad \Lini(m),\qquad \Linr(m,r).
\]
They mean, respectively, that the external invocation exists but has not
taken an abstract interval step, that its interval invocation has occurred,
and that its response has been linearized with result \(r\).

We use the paper-style assertions
\begin{align*}
  \alltok{\alpha}{\circ(m)}
    &\quad\text{all possibilities carry }\Inv(m),\\
  \sometok{\alpha}{\circ(m)}
    &\quad\text{a nonempty subconfiguration carries }\Inv(m),\\
  \alltok{\alpha}{\bullet(m)},\ \sometok{\alpha}{\bullet(m)}
    &\quad\text{similarly for }\Lini(m),\\
  \alltok{\alpha}{\bullet(m,r)},\ \sometok{\alpha}{\bullet(m,r)}
    &\quad\text{similarly for }\Linr(m,r).
\end{align*}
The connective \(P\oplusset Q\) decomposes the possibility family into two
nonempty speculative subfamilies satisfying \(P\) and \(Q\).  It is a union
connective, not ordinary propositional disjunction.  In particular, retaining
\[
  \sometok{\alpha}{\circ(\GetTop)}
  \quad\text{and}\quad
  \sometok{\alpha}{\bullet(\GetTop)}
\]
means that both fallback and snapshot explanations coexist until the concrete
return value determines which explanation is usable.

For a snapshot \(N\), the paper's named assertion is represented semantically
as
\[
\begin{split}
  \mathsf{gettingT}(\alpha,N,P)\ \triangleq\ &
  \bigl(\alltok{\alpha}{\bullet(\GetTop)}
   \land S_p(\alpha)=N\land P\bigr)
   \oplusset \mathsf{true}.
\end{split}
\]
Thus its left summand is a genuine, nonempty family of snapshot branches; the
right summand frames all alternatives not currently under discussion.

\subsection{Why rectangularity is necessary}

Different threads speculate independently about when to take their snapshot
invocations.  Merely knowing that each thread has some suitable branch is
insufficient: a later commit for \(\alpha\) must not destroy the branch needed
by a concurrent \(\beta\).

Write \(p\equiv_{\mathrm{sh}}q\) when the two pool states have identical
\(V_p,E_p,P_p,g_p\); snapshots may differ.  The invariant requires:
\begin{enumerate}
\item every pair of possibilities has shared-equivalent pool states; and
\item for any possibilities \((p_1,\pi_1),(p_2,\pi_2)\) and actor \(\alpha\),
there is a possibility \((p,\pi)\) that takes \(\alpha\)'s snapshot and token
from the first and every other thread's snapshot and token from the second:
\[
\begin{array}{ll}
 S_p(\alpha)=S_{p_1}(\alpha), & \pi(\alpha)=\pi_1(\alpha),\\
 S_p(\beta)=S_{p_2}(\beta), & \pi(\beta)=\pi_2(\beta)\quad(\beta\ne\alpha),
\end{array}
\]
with \(p\equiv_{\mathrm{sh}}p_2\).
\end{enumerate}
This Cartesian closure is the semantic counterpart of the paper's
independent-snapshot connective \(\wandsep\).  It justifies slicing the set of
possibilities by \(\alpha\)'s eventual choice while preserving every foreign
thread's local view.

\section{The global invariant}

The stable invariant has four top-level conjuncts:
\[
\boxed{\begin{aligned}
  I(\conc,\poss)\ \triangleq\ {}&
  \mathsf{ConcreteWF}(\conc)
  \land
  \bigl(\forall(\abst,\tokens)\in\poss.\
        \mathsf{BranchRep}(\conc,\abst,\tokens)\bigr)\\
  &{}\land \mathsf{Rect}(\poss)
  \land \mathsf{CounterProtocol}(\conc,\poss).
\end{aligned}}
\]
We now expand each part.

\subsection{Concrete well-formedness}

Let \(a\) and \((\tau,q)\) be the array and timestamp components of
\(\conc\).  Then \(\mathsf{ConcreteWF}\) contains:
\begin{align*}
&q(\alpha)=l \Longrightarrow l\le\tau,
\tag{WF-TS}\\
&T(n)=\interval{l}{u}\Longrightarrow u+1\le\tau,
\tag{WF-Clock}\\
&V(n)=v\Longrightarrow
  \exists t.\ T(n)=t\land\mathsf{wf}_{\TS}(t),
\tag{WF-Domain-1}\\
&T(n)=t\Longrightarrow\exists v.\ V(n)=v,
\qquad
 n\in g\Longrightarrow n\in\dom(V),
\tag{WF-Domain-2}\\
&n=(\beta,\ell)\text{ is live}
  \Longleftrightarrow \ell\in O(\beta),
\tag{WF-Order}\\
&n\in\dom(V)\Longrightarrow\beta\in\Upsilon,
\qquad \mathsf{NoDup}(O(\beta)).
\tag{WF-Threads}
\end{align*}
The clock property is the bridge between completed timestamps and future
timestamp invocations.  The order properties connect each row-level
\(\mathsf{getTop}\) response to the graph argument.

\subsection{Branch representation}

Every stable abstract possibility is \(\Ready(p)\), and
\(\mathsf{BranchRep}(\conc,\Ready(p),\tokens)\) consists of three pieces.

First, \(p\) represents the concrete array:
\begin{align*}
V_p(n)&=V(n),
\tag{Rep-V}\\
E_p(n_1,n_2)&\Longrightarrow \arrayedge(n_1,n_2),
\tag{Rep-E}\\
E_p(n_1,n_2)&\Longrightarrow n_1,n_2\in\dom(V_p),
\tag{Rep-Dom}\\
n\in g_p&\Longleftrightarrow n\in g,
\tag{Rep-G}\\
n\in |P_p|&\Longleftrightarrow T(n)=\topts,
\tag{Rep-P}\\
S_p(\alpha)=N\land n\in N&\Longrightarrow n\in\dom(V).
\tag{Rep-S}
\end{align*}
Within one row, abstract edges between live vertices agree with allocation
order, and any two distinct live vertices in that row are comparable by an
abstract edge.  This is stronger than simple edge inclusion and is exactly
what identifies the row response as the row's abstract top.

Second, the pool control state and token map obey a protocol:
\begin{align*}
P_p(\alpha)=\ell
  &\Longrightarrow \exists v.\ \pi(\alpha)=\Lini(\Push(v)),\\
S_p(\alpha)=N
  &\Longrightarrow \pi(\alpha)=\Lini(\GetTop),\\
\pi(\alpha)=\Lini(\Push(v))
  &\Longrightarrow \exists\ell.\ P_p(\alpha)=\ell.
\end{align*}

Third, pending timestamp invocations retain a causal edge fact.  Define
\[
\mathsf{outgoingBefore}(a,p,n,l)\ \triangleq\
  \forall m.\ E_p(n,m)\Longrightarrow
  \exists l_m,u_m.\ T(m)=\interval{l_m}{u_m}\land u_m<l.
\]
Then
\begin{align*}
q(\alpha)=l&\Longrightarrow\exists\ell.\ P_p(\alpha)=\ell,\\
q(\alpha)=l\land P_p(\alpha)=\ell
  &\Longrightarrow
  \mathsf{outgoingBefore}(a,p,(\alpha,\ell),l).
\tag{Rep-Causal}
\end{align*}
When \(\mathsf{newTS}\) returns \(\interval{l}{u}\), every abstract outgoing
edge of that push therefore becomes a concrete timestamp edge.  This proves
\(E_p\subseteq\arrayedge\) after \(\mathsf{setTS}\).

\subsection{Pending aggregate counters}

The final invariant clause states
\[
PC(\alpha)=c
\quad\Longrightarrow\quad
\exists(\rho,\pi)\in\Delta.\ \pi(\alpha)=\Lini(\GetTop).
\]
An internal aggregate-counter call can therefore be pending only when an
abstract \(\GetTop\) interval invocation exists.  This rules out stale
counter state at a fresh external invocation and is essential for both
stability and error reasoning.

\section{Speculative witnesses used by the scan}

\subsection{Candidate and row witnesses}

Suppose caller \(\alpha\) has completed rows in list \(A\), and the
implementation accumulator contains
\[
  r=(v,\beta,\ell,t),\qquad n_r=(\beta,\ell).
\]
A \emph{candidate view} is an actual possibility
\((\Ready(p),\pi)\in\Delta\) with a snapshot \(S_p(\alpha)=N\) such that
\begin{align*}
&\pi(\alpha)=\Lini(\GetTop),\quad V_p(n_r)=v,\quad \beta\in A,\\
&n_r\in g_p\quad\text{or}\quad
  \Bigl(n_r\in N\setminus g_p\ \land
   \forall n'\in N\setminus g_p.\
     \fst(n')\in A\Rightarrow\neg E_p(n',n_r)\Bigr).
\tag{Cand-Top}
\end{align*}
If the observed timestamp was \(\topts\), the view additionally asserts
\[
  \forall n'\in N\setminus g_p.\ \neg E_p(n',n_r).
\tag{Cand-TopTS}
\]
The stronger condition is needed because the timestamp saved in \(r\) does
not change when the node is later stamped.

During a pending row call, a \emph{row snapshot view} ties the saved concrete
row order \(L\) to the same abstract possibility: all addresses in \(L\) are
in \(N\); every live snapshot node of that row occurs in \(L\); and abstract
edges between live snapshot nodes respect the order in \(L\).  The proof uses
a combined candidate--row witness rather than conjoining two existential
assertions.  This prevents accidentally taking the candidate from one
speculative branch and row coverage from an unrelated branch.

\subsection{The node-cut invariant}

The paper motivates deletion of ignored vertices from a snapshot.  A literal
induction over ignored nodes would require finite support for the
predicate-valued historical vertex set, which the state representation does
not provide.  The mechanized proof instead maintains the exact consequence
needed at commit time.

For a live, nonignored node \(n\), a \emph{node cut} is a real possibility
with snapshot \(N\) such that
\begin{equation}
\begin{split}
&\pi(\alpha)=\Lini(\GetTop),\quad S_p(\alpha)=N,
  \quad n\in N,\quad V_p(n)=V(n),\\
&\forall n'\in N\setminus g_p.\
  E_p(n',n)\Longrightarrow
  \mathsf{status}_{\alpha}(n')\ne\mathsf{Ignored}.
\end{split}
\tag{Cut}
\end{equation}
The global scan assertion requires such a cut for every concrete live node
that is not ignored by the current scan.  Old cuts survive a push that is
ignored by the scan; a newly inserted nonignored node receives a fresh cut in
the newly forked snapshot family.  At the end of a complete scan, every
nonignored live predecessor has been visited.  Candidate maximality excludes
all such predecessors, so the cut witnesses that the candidate is a genuine
top.  This proof-relevant construction replaces the paper's finite repeated
removal argument without adding a finiteness axiom.

\subsection{The stale-\texorpdfstring{\(\topts\)}{TSTop} problem}

There is a tempting but false loop invariant: ``the saved candidate is a top
under the current concrete timestamp order.''  A scan may read \(n\) while
\(T(n)=\topts\), after which the owning push replaces \(\topts\) by an interval.
The saved candidate still contains \(\topts\), so its comparison behavior no
longer matches the current array state.

The proof separates two cases:
\begin{itemize}
\item If the saved timestamp is an interval, it is immutable.  The assertion
  records that the current concrete timestamp is still that exact interval.
\item If the saved timestamp is \(\topts\), no equality with the current
  timestamp is required.  Instead, condition \textup{(Cand-TopTS)} is retained.
\end{itemize}
Why is \textup{(Cand-TopTS)} stable?  Any push that begins while \(n\) is still
pending sees \(n\in|P_p|\), and the abstract push rule deliberately does not
add an edge to a pending node.  Completing \(n\)'s timestamp later changes the
concrete timestamp order but adds no abstract edge.  Older completed pushes
cannot become incoming predecessors of \(n\), and later pushes not present in
the saved snapshot are irrelevant.  This causal argument covers precisely the
case that timestamp comparison alone cannot handle.

\section{Rely and guarantee relations}

The Hoare judgment has the form
\[
  \{P\}_{R_\alpha,G_\alpha,I}\ c\ \{x.Q(x)\},
\]
meaning that thread \(\alpha\) executes \(c\), environment steps satisfy
\(R_\alpha\), implementation steps satisfy \(G_\alpha\), and all stable
boundaries satisfy \(I\).

\subsection{Evolution relations}

Several monotone relations summarize exactly what interference may do.

\paragraph{Array evolution.}
For every owner \(\beta\),
\[
C(\beta)\le C'(\beta),
\qquad
C(\beta)=C'(\beta)\Longrightarrow O'(\beta)\subseteq O(\beta).
\tag{Arr-Evol}
\]
Thus a counter can increase only because of insertion; with no increase, the
live order can only shrink by removal.

\paragraph{Garbage and intervals.}
Garbage is monotone, timestamp domains do not shrink, and an existing interval
timestamp is immutable:
\begin{align*}
g&\subseteq g',\\
T(n)\downarrow&\Longrightarrow T'(n)\downarrow,\\
T(n)=\interval{l}{u}&\Longrightarrow T'(n)=\interval{l}{u}.
\end{align*}
The timestamp clock is also monotone.

\paragraph{Local concrete state.}
For observer \(\beta\), its scan progress and pending aggregate-counter entry
are unchanged by foreign steps.  Its pending timestamp lower bound is likewise
unchanged.

\paragraph{Token evolution.}
Old token views are retained.  The only new view permitted for observer
\(\beta\) is \(\Lini(\GetTop)\), and it must be justified by an old
\(\Inv(\GetTop)\) view.  Formally,
\begin{align*}
\mathsf{view}_\beta(\Delta,t)&\Rightarrow
  \mathsf{view}_\beta(\Delta',t),\\
\mathsf{view}_\beta(\Delta',t)&\Rightarrow
  \mathsf{view}_\beta(\Delta,t)\ \lor\
  \bigl(t=\Lini(\GetTop)\land
        \mathsf{view}_\beta(\Delta,\Inv(\GetTop))\bigr).
\end{align*}
This is the logical permission to add snapshot branches while retaining the
atomic fallback.

\paragraph{Abstract local views.}
Existing pairs consisting of an observer's pending-push entry and snapshot
entry remain available.  New snapshot alternatives may appear, but the
pending-push component of every new view agrees with some old view.  Candidate
views, row views, combined candidate--row views, and node cuts are preserved.
The causal fact \(\mathsf{outgoingBefore}\) for the observer's pending push is
also preserved.

\subsection{Definitions of \texorpdfstring{\(G_\alpha\) and \(R_\alpha\)}{G and R}}

The guarantee \(G_\alpha(w,w')\) is the conjunction of all evolution clauses
above, with observer-indexed clauses required for every
\(\beta\ne\alpha\).  The rely \(R_\alpha(w,w')\) is the same conjunction
specialized to observer \(\alpha\).  Schematically,
\[
\begin{split}
G_\alpha={}&
 \bigwedge_{\beta\ne\alpha}
  (\mathsf{TokenRely}_\beta\land\mathsf{PoolLocal}_\beta
   \land\mathsf{LocalConcrete}_\beta
   \land\mathsf{WitnessPreservation}_\beta)
\\&{}land\mathsf{ArrEvol}\land\mathsf{ClockMono}
 \land\mathsf{GarbageEvol}\land\mathsf{IntervalEvol},\\
R_\alpha={}&
 \mathsf{TokenRely}_\alpha\land\mathsf{PoolLocal}_\alpha
 \land\mathsf{LocalConcrete}_\alpha
 \land\mathsf{WitnessPreservation}_\alpha
\\&{}land\mathsf{ArrEvol}\land\mathsf{ClockMono}
 \land\mathsf{GarbageEvol}\land\mathsf{IntervalEvol}.
\end{split}
\]

The construction immediately suggests the compatibility proof:
\[
  \alpha\ne\beta\quad\Longrightarrow\quad G_\alpha\subseteq R_\beta.
\]
Identity, invocation, and response bookkeeping steps satisfy the corresponding
relies, and both relations are transitive.  Hence the pair \((R,G)\) is a valid
parallel rely--guarantee assignment.  The long preservation arguments for
push, removal, snapshot forks, and possibility restriction exist to establish
the witness-preservation clauses of this inclusion.

\section{Generic method assertions and Hoare rules}

For any exported operation \(m\), define
\[
\begin{aligned}
\mathsf{Active}_\alpha(m)&\triangleq
 I\land\alltok{\alpha}{\circ(m)},\\
\mathsf{Completed}_\alpha(m,r)&\triangleq
 I\land\alltok{\alpha}{\bullet(m,r)}.
\end{aligned}
\]
Both assertions are stable under \(R_\alpha\), except that an active
\(\GetTop\) needs a specialized existential entry assertion because
interference is allowed to add snapshot branches.

The proof repeatedly uses three semantic Hoare principles.

\begin{description}
\item[Underlay call.] To prove a visible call, show that its invocation cannot
error, prove an invocation update \(P\tran P'\), prove one response update
\(P'\tran Q(r)\) for each return \(r\), and establish stability of the
intermediate assertions.

\item[Error split.] Before an underlay call, refine \(P\) into either a safe
precondition or an abstract-error possibility.  This handles calls outside
the finite domain and undefined removals without assuming safety.

\item[Logical update.] A concrete-free update may restrict or extend the
possibility family along zero or more legal abstract steps.  It is used to
commit to a candidate slice and to add the counter-result alternatives.
\end{description}

The crucial discipline is that ordinary deterministic underlay steps use a
pointwise image of \(\Delta\), preserving one descendant for every prior
possibility.  The proof restricts the family only when the concrete observation
has supplied enough information to justify a commit.

\section{The \texorpdfstring{\(\TryRemove\)}{tryRemove} triple}

The target triple is
\[
\boxed{
\{\mathsf{Active}_\alpha(\TryRemove(\beta,\ell))\}
\ \mathsf{array.tryRemove}(\beta,\ell)\ 
\{b.\mathsf{Completed}_\alpha(\TryRemove(\beta,\ell),b)\}.
}
\]

The precondition is split into a safe array call or an abstract error.  If
\(\alpha\notin\Upsilon\), if \(\beta\notin\Upsilon\), or if the node is
undefined, branch representation supplies the corresponding ListPool error
step.  Otherwise the array invocation is safe.

On a successful response, the concrete array adds
\(n=(\beta,\ell)\) to \(g\).  Every possibility performs the matching atomic
ListPool invocation and successful response, adds \(n\) to \(g_p\), and changes
\(\alpha\)'s token to \(\Linr(\TryRemove(n),\true)\).  On failure, branch
representation gives \(n\in g_p\); every possibility performs the failed
abstract response and leaves the graph unchanged.  Because garbage is
monotone, an existing candidate either remains live/top or moves to the
garbage alternative.  Row views and node cuts survive: weakening liveness by
adding garbage cannot create a new live predecessor.  This establishes the
guarantee and the postcondition for both booleans.

\section{The \texorpdfstring{\(\Push\)}{push} triple}

The target triple is
\[
\boxed{
\{\mathsf{Active}_\alpha(\Push(v))\}
\ \ell\leftarrow\mathsf{insert}(v);\,
 t\leftarrow\mathsf{newTS}();\,
 \mathsf{setTS}(\ell,t)\ 
\{\_.\mathsf{Completed}_\alpha(\Push(v),())\}.
}
\]

The useful intermediate assertions are
\begin{align*}
\mathsf{PushActor}(\alpha,v)
  &\triangleq \mathsf{Active}_\alpha(\Push(v))\land\alpha\in\Upsilon,
\\
\mathsf{PushLin}(\alpha,v,\ell)
  &\triangleq I\land\alltok{\alpha}{\bullet(\Push(v))}
    \land\alpha\in\Upsilon
    \land \forall\text{ local views}.\ P_p(\alpha)=\ell,
\\
\mathsf{PushStamped}(\alpha,v,\ell,l,u)
  &\triangleq \mathsf{PushLin}(\alpha,v,\ell)
    \land l\le u\land u+1\le\tau
\\&\qquad{}land q(\alpha)\uparrow
    \land\mathsf{outgoingBefore}(a,p,(\alpha,\ell),l).
\end{align*}

\subsection{Insertion and speculative snapshot saturation}

The array insertion allocates \(n=(\alpha,\ell)\) with timestamp \(\topts\),
increments its row counter, and adds it to the row head.  Every abstract branch
takes the ListPool push invocation, adding the same value and the appropriate
outgoing edges.  Thus
\[
  \mathsf{PushActor}(\alpha,v)
  \tran \mathsf{PushLin}(\alpha,v,\ell).
\]

Insertion may affect every currently scanning thread: if the new node lies in
an unvisited row, some future snapshot explanation should contain it; if its
row was already snapshotted, it is ignored for that scan.  The possibility
family is therefore saturated by optional snapshot forks for the participating
threads.  Each fork retains the old branch and, when legal, adds the branch
that takes a ListPool snapshot invocation.  Repeated optional forks preserve
rectangularity and all old token/local views.  The node-cut invariant is
updated constructively: ignored insertions preserve old cuts, whereas a
nonignored insertion obtains a fresh cut from a new snapshot branch.

\subsection{Timestamp invocation and response}

Invocation of \(\mathsf{newTS}\) records the current clock lower bound \(l\).
The branch invariant implies that every outgoing abstract predecessor of the
pending push already carries an interval whose upper endpoint is below \(l\).
This is exactly \(\mathsf{outgoingBefore}\).

The timestamp response returns only an interval \(\interval{l}{u}\) with
\(l\le u\), clears \(q(\alpha)\), and advances the clock beyond \(u\).  A
sentinel response is impossible and is assigned the false assertion.  Hence
the response establishes \(\mathsf{PushStamped}\).

\subsection{Installing the timestamp and completing the push}

The final array call replaces \(T(n)=\topts\) by
\(T(n)=\interval{l}{u}\).  For every abstract edge \(E_p(n,m)\), causal
ordering gives \(T(m)=\interval{l_m}{u_m}\) with \(u_m<l\); therefore
\[
  \interval{l_m}{u_m}<_{\TS}\interval{l}{u}
  \quad\text{and hence}\quad \arrayedge(n,m).
\]
All old interval timestamps are unchanged.  Every possibility takes the
abstract push response, removes \(\alpha\) from \(P_p\), and sets its token to
\(\Linr(\Push(v),())\).  The representation and protocol invariants are
restored, yielding \(\mathsf{Completed}_\alpha(\Push(v),())\).

\section{The \texorpdfstring{\(\GetTop\)}{getTop} triple}

This is the core of the proof.  Its external triple is
\[
\boxed{
\{\mathsf{GetTopEntry}(\alpha)\}
\ \mathsf{getTopImpl}_\alpha\ 
\{r.\mathsf{Completed}_\alpha(\GetTop,r)\}.
}
\]
The entry assertion contains \(I\), an existential unlinearized fallback
\(\sometok{\alpha}{\circ(\GetTop)}\), and either a valid in-domain local state
or evidence that \(\alpha\notin\Upsilon\).  The latter case takes the abstract
error step.  In the former case the pending-counter protocol proves that
\(PC(\alpha)\) is absent.

\subsection{Reset and the initial fork}

The array reset installs empty scan progress for \(\alpha\).  At its response,
the proof applies an optional abstract snapshot step:
\[
  \Delta' = \Delta
   \ \cup\ 
   \{\text{descendants obtained by taking }\alpha\text{'s snapshot invocation}\}.
\]
Thus the reset postcondition contains both
\[
  \sometok{\alpha}{\circ(\GetTop)}
  \quad\text{and}\quad
  \sometok{\alpha}{\bullet(\GetTop)},
\]
at least one concrete snapshot view, and node cuts for all initially live
nodes.  The first token family is reserved for atomic \(\Fail\),
\(\Succ(\Empty)\), and garbage explanations; the second supplies nonempty
top explanations.

\subsection{Loop invariant}

Let the fixed thread list decompose as
\[
  \Upsilon=A\mathbin{++}B,
\]
where \(A\) has been scanned and \(B\) remains.  Let the concrete accumulator
be \((r,c)\), with optional candidate \(r\) and empty-row count \(c\).  The
stable assertion \(\mathsf{ScanFold}_\alpha(B,r,c)\) contains:

\begin{enumerate}
\item the global invariant \(I\), domain membership of \(\alpha\), absence of
its pending aggregate counter, and exact scan progress with visited rows
\(A\), no current row, and all seen nodes owned by rows in \(A\);

\item both token families
\(\sometok{\alpha}{\circ(\GetTop)}\) and
\(\sometok{\alpha}{\bullet(\GetTop)}\), a snapshot possibility, and available
node cuts;

\item every seen node has a timestamp in the concrete domain; and

\item the following accumulator condition.
\end{enumerate}

If \(r=\mathsf{None}\), define
\[
\mathsf{EmptyEvidence}(A,c,a)\ \triangleq\
c\le\sum_{\beta\in A}C(\beta)
\ \land\ 
\left(c=\sum_{\beta\in A}C(\beta)
 \Rightarrow \forall\beta\in A.\ O(\beta)=[]\right).
\tag{Empty-Ev}
\]
The invariant requires \(\mathsf{EmptyEvidence}\) and says every node seen so
far is garbage.  This assertion is stable by \textup{(Arr-Evol)}: if the total
does not increase, each component counter is unchanged and each current row
order is included in the previously empty order.

If \(r=\mathsf{Some}(v,\beta,\ell,t)\), the invariant contains a candidate
view, timestamp validity, and maximality over seen live nodes:
\[
  \forall n,t'.\ \mathsf{seen}(n)\land n\notin g\land T(n)=t'
  \Longrightarrow \neg(t<_{\TS}t').
\tag{Cand-Max}
\]
For interval candidates, timestamp validity is exact equality with the current
array timestamp.  For \(\topts\) candidates, the separate causal safety
condition described earlier is used.

During the call on the next row \(\beta\), a pending assertion additionally
records the saved row order and counter, a row snapshot view, no duplicates,
and the combined candidate--row witness.  It also records
\[
 C_{\mathrm{saved}}(\beta)\le C(\beta),
\qquad
C_{\mathrm{saved}}(\beta)=C(\beta)
 \Rightarrow O(\beta)\subseteq O_{\mathrm{saved}}(\beta).
\]

\subsection{One row-step triple}

For \(B=\beta::B'\), the loop body satisfies
\[
\boxed{
\{\mathsf{ScanFold}_\alpha(\beta::B',r,c)\}
\ \mathsf{scanStep}(\beta,r,c)\ 
\{(r',c').\mathsf{ScanFold}_\alpha(B',r',c')\}.
}
\]

At invocation of \(\mathsf{array.getTop}(\beta)\), the saved row order is
linked to a snapshot possibility and the scan state changes from idle to
visiting.  The response has two cases.

\paragraph{Nonempty row.}
Suppose the row returns \((v,t,\ell)\).  The saved-order semantics makes
\((\beta,\ell)\) the first currently live address in the saved order.  Row
ordering then excludes any live incoming predecessor from that row in the
same snapshot.  The returned node receives a candidate view, using its node
cut when necessary.  The implementation chooses between it and the old
candidate by the boolean version of \(<_{\TS}\).  In the ``newer'' branch,
the new timestamp is strictly above the old saved interval and becomes the
maximal candidate.  In the retained branch, the failed comparison directly
preserves \textup{(Cand-Max)}.  All combinations involving \(\topts\) are
handled by the causal safety predicate, so no equality between a stale saved
\(\topts\) and the current timestamp is assumed.

\paragraph{Empty row.}
Suppose the row returns saved counter \(d\).  The implementation changes
\(c\) to \(c+d\).  If no candidate exists, the row response proves the new
empty-evidence bound.  If equality with the current row counter holds, the
current live order is empty; otherwise the inequality records possible
interference.  If a candidate already exists, an empty response cannot
introduce a larger seen live node, so its candidate view and maximality are
preserved.

In both cases the current row is appended to \(A\), its snapshot nodes become
seen, and the scan returns to the idle phase.  The generic finite-list
\(\mathsf{foreach}\) rule iterates this triple to obtain
\[
  \mathsf{ScanFold}_\alpha(\Upsilon,\mathsf{None},0)
  \tran \exists r,c.\mathsf{ScanFold}_\alpha([],r,c).
\]

\subsection{Committing a nonempty candidate}

At loop exit with candidate \(r=(v,\beta,\ell,t)\), the candidate view supplies
a donor possibility \((p_d,\pi_d)\) with snapshot \(N\).  Since all rows have
been visited, condition \textup{(Cand-Top)} strengthens to one of:
\[
  n_r\in g_{p_d}
  \qquad\text{or}\qquad
  \mathsf{top}(N\setminus g_{p_d},E_{p_d},n_r).
\]
In the live case, the node-cut invariant excludes ignored incoming
predecessors and candidate maximality excludes nonignored ones.  In the
garbage case, the relaxed abstract garbage-response rule applies directly.

The commit has two mathematical stages.

\paragraph{Actor slice.}
Restrict \(\Delta\) to branches whose \(\alpha\)-snapshot and \(\alpha\)-token
equal the donor's choices:
\[
  \mathsf{Slice}_\alpha(\Delta,p_d,\pi_d).
\]
Rectangularity proves this slice nonempty and, for every foreign observer,
still representative of all its former local choices.  This is the formal
meaning of selecting one \(\oplusset\)-summand without invalidating concurrent
proofs.

\paragraph{Pointwise response image.}
Every branch in the slice takes either the abstract top response or the
abstract garbage response, clears \(S_p(\alpha)\), and changes the token to
\(\Linr(\GetTop,\Succ(v,\beta,\ell))\).  Applying this response pointwise
produces \(\Delta'\).  Shared graph equality transfers the donor's top/garbage
fact to every branch in the slice.  The result is
\[
\mathsf{ScanFold}_\alpha([],\mathsf{Some}(r),c)
\tran
\mathsf{Completed}_\alpha(\GetTop,\Succ(v,\beta,\ell)).
\]

\subsection{The empty/failure counter protocol}

At loop exit with no candidate, let the accumulated count be \(c\).  From
\textup{(Empty-Ev)} over all rows,
\[
  c\le C_{\mathrm{tot}}
  \quad\text{and}\quad
  c=C_{\mathrm{tot}}\Rightarrow
  \forall\beta\in\Upsilon.\ O(\beta)=[].
\tag{All-Empty}
\]

Invocation of \(\mathsf{array.getCounter}\) saves
\(s=C_{\mathrm{tot}}\) and later returns some \(z\) satisfying
\[
  s\le z\le C'_{\mathrm{tot}}.
\tag{Counter-Bounds}
\]
At invocation time, the proof extends the possibility family with an atomic
\(\Fail\) response alternative unconditionally.  If \(c=s\),
\textup{(All-Empty)} and the live-order representation imply that every
abstract vertex is garbage, so it also adds an atomic
\(\Succ(\Empty)\) alternative.  Old unlinearized and snapshot branches are
retained.

At response time:
\begin{itemize}
\item If \(z=c\), then \(c\le s\le z=c\), hence \(c=s\).  The empty
alternative was created and is legal; slice to it and return
\(\Succ(\Empty)\).
\item If \(z\ne c\), slice to the always-available failure alternative and
return \(\Fail\).
\end{itemize}
This is exactly the concrete branch
\[
  \mathsf{if}\ z=c\ \mathsf{then}\ \Succ(\Empty)\ \mathsf{else}\ \Fail.
\]
The counter argument does not claim that a failed scan saw a single atomic
empty instant.  It only claims emptiness when equality rules out every
overlapping insertion; otherwise the permissive abstract failure transition
accounts for the execution.

\section{Assembly of the Hoare proof}

The three method triples are now:
\begin{align*}
\{\mathsf{Active}_\alpha(\Push(v))\}\;&\mathsf{pushImpl}(v)\;
 \{r.\mathsf{Completed}_\alpha(\Push(v),r)\},\\
\{\mathsf{GetTopEntry}(\alpha)\}\;&\mathsf{getTopImpl}\;
 \{r.\mathsf{Completed}_\alpha(\GetTop,r)\},\\
\{\mathsf{Active}_\alpha(\TryRemove(n))\}\;&\mathsf{tryRemoveImpl}(n)\;
 \{r.\mathsf{Completed}_\alpha(\TryRemove(n),r)\}.
\end{align*}

For \(\GetTop\), sequential composition has the following proof outline:
\[
\begin{array}{c}
\mathsf{GetTopEntry}
\\ \Downarrow\ \mathsf{resetIter}
\\ \mathsf{GetTopReset}
\\ \Downarrow\ \text{logical initialization}
\\ \mathsf{ScanFold}(\Upsilon,\mathsf{None},0)
\\ \Downarrow\ \mathsf{foreach}\text{ using the row triple}
\\ \exists r,c.\ \mathsf{ScanFold}([],r,c)
\\ \Downarrow
\\
\begin{cases}
\mathsf{candidate\ slice+response},&r=\mathsf{Some}(r_0),\\
\mathsf{counter\ invocation+response},&r=\mathsf{None},
\end{cases}
\\ \Downarrow
\\ \mathsf{Completed}_\alpha(\GetTop,\mathsf{result}).
\end{array}
\]

The initial concrete state satisfies \(I\): maps and garbage are empty,
counters and row orders are initialized for \(\Upsilon\), the timestamp clock
is zero, and the abstract configuration is the singleton empty pool with an
empty token map.  Singleton configurations are rectangular.  The
rely--guarantee assignment is valid and parallel-compatible, and each method's
pre/post assertions expose and discharge the generic invocation/response
tokens.  Soundness of the set-valued rely--guarantee logic therefore yields
the layer simulation.

\begin{theorem}[ListPool layer correctness]
For every value type and finite thread domain \(\Upsilon\), the implementation
of ListPool using \(\mathsf{SPListArray}(\Upsilon)\otimes\mathsf{Timestamp}\)
simulates the interval-sequential DAG specification of ListPool.  Consequently
the implementation layer is contextually linearizable with respect to that
specification.
\end{theorem}

\clearpage
\section{Coq proof statistics}

\subsection{Whole-development footprint}

The following table reports the size of the primary Coq development for this
layer.  ``Physical'' is the result of \code{wc -l}.  The remaining three
columns use \code{coqwc}'s classification: declaration and term lines are
reported as ``spec'', tactic-script lines as ``proof'', and comment-only lines
as ``comments''.  Blank and layout-only lines explain why the \code{coqwc}
columns do not sum to the physical count.

\begin{center}
\centering
\begin{tabular}{lrrrr}
\toprule
File & Physical & \code{coqwc} spec & \code{coqwc} proof & Comments \\
\midrule
\code{ListPool.v}       &    166 &   118 &    14 & 10 \\
\code{ListPoolSpec.v}   &    355 &   269 &    22 & 21 \\
\code{TimestampSpec.v}  &    202 &   122 &    37 & 13 \\
\code{ListPoolProof.v}  & 10,953 & 2,854 & 7,534 & 92 \\
\midrule
Total                    & 11,676 & 3,363 & 7,607 & 136 \\
\bottomrule
\end{tabular}
\end{center}

Within \code{ListPoolProof.v}, the 2,854 declaration/term lines, 7,534 proof
lines, and 92 comment lines leave 473 blank or layout-only physical lines.
Thus tactic script accounts for 68.8\% of the physical file, and the ratio of
classified proof lines to classified declaration lines is 2.64.  The three
supporting files add only 73 proof lines; 99.0\% of the 7,607 proof lines in
the four-file development are in \code{ListPoolProof.v} itself.

\subsection{Method-level footprint}

The next table partitions the proof file into its shared foundation, the
three method-specific regions, and final packaging.  These are syntactic,
non-overlapping regions: the method rows do not charge reused foundational
lemmas to each consumer.

\begin{center}
\small
\begin{tabular}{lrrrrrr}
\toprule
Region & Source lines & Physical & Spec & Proof & \% physical & \% proof \\
\midrule
Shared foundation & 1--3,408 & 3,408 & 1,172 & 1,959 & 31.1 & 26.0 \\
\code{tryRemove} & 3,409--4,031 & 623 & 57 & 556 & 5.7 & 7.4 \\
\code{push} & 4,032--5,709 & 1,678 & 248 & 1,378 & 15.3 & 18.3 \\
\code{getTop} & 5,710--10,903 & 5,194 & 1,369 & 3,601 & 47.4 & 47.8 \\
Layer packaging & 10,904--10,953 & 50 & 8 & 40 & 0.5 & 0.5 \\
\midrule
Total & 1--10,953 & 10,953 & 2,854 & 7,534 & 100.0 & 100.0 \\
\bottomrule
\end{tabular}
\end{center}

Nearly half of both the physical source and the classified proof script is
dedicated to \code{getTop}.  Its 3,601 proof lines are 2.61 times the size of
the \code{push} proof and 6.48 times the size of the \code{tryRemove} proof.
This is consistent with the mathematical structure: \code{getTop} alone needs
speculative snapshot families, the row fold, candidate commitment, and the
empty/failure counter protocol.

\subsection{Phase-by-phase profile}

A finer architectural partition is shown below.  ``Defs'' combines ordinary
and program definitions.  The \code{coqwc} columns were computed on each
source interval with header suppression disabled; comments and layout lines
are omitted from this table, so Physical need not equal Spec plus Proof.

\begin{center}
\scriptsize
\begin{tabular}{lrrrrrr}
\toprule
Phase & Source lines & Physical & Spec & Proof & Lemmas & Defs \\
\midrule
Setup and global invariant & 1--469 & 469 & 247 & 153 & 14 & 25 \\
Possibilities and rely--guarantee & 470--1,595 & 1,126 & 464 & 549 & 49 & 29 \\
Representation preservation and removal & 1,596--4,031 & 2,436 & 518 & 1,813 & 88 & 6 \\
Push triple & 4,032--5,709 & 1,678 & 248 & 1,378 & 36 & 12 \\
\code{getTop} snapshot and reset & 5,710--6,623 & 914 & 252 & 608 & 39 & 10 \\
\code{getTop} loop and row calls & 6,624--8,817 & 2,194 & 548 & 1,565 & 52 & 15 \\
\code{getTop} candidate commit & 8,818--9,521 & 704 & 198 & 473 & 24 & 3 \\
\code{getTop} empty/failure counter & 9,522--10,655 & 1,134 & 277 & 814 & 37 & 6 \\
Triple assembly and layer theorem & 10,656--10,953 & 298 & 102 & 181 & 12 & 3 \\
\midrule
Total & 1--10,953 & 10,953 & 2,854 & 7,534 & 351 & 109 \\
\bottomrule
\end{tabular}
\end{center}

The loop-and-row phase is the largest \code{getTop} subproof, with 1,565
classified proof lines.  The counter phase contributes another 814 lines, and
the candidate-slicing commit contributes 473.  The large
``representation preservation and removal'' phase includes the common
preservation library used by all methods as well as the dedicated
\code{tryRemove} triple; it is therefore intentionally broader than the
method-level \code{tryRemove} row above.

\subsection{Declaration and obligation profile}

The proof file contains 351 \code{Lemma} declarations, 104 ordinary
\code{Definition}s, five \code{Program Definition}s, five auxiliary
\code{Variant}s, and two recursive \code{Fixpoint}s.  The five program
definitions generate nine explicit \code{Next Obligation} proofs.  Hence the
file has 360 completed \code{Qed} proof blocks: 351 named lemmas and nine
generated obligations.  The definitions are mostly logical assertions,
possibility transformers, and auxiliary relations, not executable ListPool
components.

For a useful measure of proof granularity, define a lemma's \emph{declaration
span} as the inclusive physical-line interval from its \code{Lemma} command to
the line containing its closing \code{Qed}.  This includes the statement,
proof script, comments, and local whitespace.  The 351 spans total 9,355 lines,
have mean 26.7, median 16, minimum 2, and maximum 493 lines.

\begin{center}
\begin{tabular}{lrr}
\toprule
Declaration span & Lemmas & Percentage \\
\midrule
2--10 lines & 117 & 33.3\% \\
11--25 lines & 129 & 36.8\% \\
26--50 lines & 70 & 19.9\% \\
51--100 lines & 22 & 6.3\% \\
101--200 lines & 8 & 2.3\% \\
More than 200 lines & 5 & 1.4\% \\
\midrule
Total & 351 & 100.0\% \\
\bottomrule
\end{tabular}
\end{center}

Thus 70.1\% of lemmas occupy at most 25 lines, while the small number of large
response-update proofs accounts for a substantial fraction of the file.  The
ten largest lemma spans are:

\begin{center}
\footnotesize
\begin{tabular}{>{\raggedright\arraybackslash}p{0.42\textwidth}rr
                >{\raggedright\arraybackslash}p{0.30\textwidth}}
\toprule
Lemma & Source lines & Span & Main obligation \\
\midrule
\code{push_insert_res_update} & 4,390--4,882 & 493 & Insert response and snapshot saturation \\
\code{getTop_row_inv_update} & 7,233--7,521 & 289 & Begin a row scan and align witnesses \\
\code{tryRemove_true_res_update} & 3,505--3,754 & 250 & Successful removal response \\
\code{push_setTS_res_update} & 5,270--5,519 & 250 & Install timestamp and finish push \\
\code{getTop_row_empty_res_update} & 8,310--8,549 & 240 & Empty-row response and counter evidence \\
\code{pool_represents_start_push} & 1,771--1,970 & 200 & Preserve the graph representation at push \\
\code{getTop_reset_res_update} & 6,424--6,622 & 199 & Reset response and optional snapshot fork \\
\code{getTop_row_nonempty_res_update} & 8,551--8,748 & 198 & Nonempty-row candidate update \\
\code{tryRemove_false_res_update} & 3,756--3,948 & 193 & Failed removal response \\
\path{possibility_rectangular_add_actor_alternative} & 9,795--9,915 & 121 & Preserve rectangularity for counter alternatives \\
\bottomrule
\end{tabular}
\end{center}

Together these ten declarations span 2,433 lines, or 26.0\% of all lemma
spans.  Their concentration around response updates reflects the central proof
burden: each observed underlay response must update every abstract possibility,
restore \(I\), and establish all foreign-thread guarantee clauses.

The development contains no \code{Admitted}, \code{admit}, \code{Axiom}, or
\code{Abort} escape hatch.

\subsection{Reproducing the counts}

The counts can be reproduced from the repository root with
\begin{quote}
\small
\begin{verbatim}
wc -l examples/TSStack/ListPool*.v \
  examples/TSStack/TimestampSpec.v
coqwc examples/TSStack/ListPool*.v \
  examples/TSStack/TimestampSpec.v
\end{verbatim}
\end{quote}

\section{Where the mechanization strengthens the paper argument}

Three points are worth making explicit for readers comparing this proof with
the paper outline.

\begin{enumerate}
\item The proof uses a genuinely set-valued abstract configuration throughout.
The atomic fallback, interval-snapshot branches, and independent choices for
different scans coexist.  No singleton abstraction can express the required
future-dependent commit.

\item Rectangularity is maintained as an invariant, not treated informally.
It is what makes actor-local restriction a valid guarantee step for every
other thread.

\item The node-cut witness replaces induction over an implicitly finite set of
ignored historical vertices.  This both avoids a hidden finiteness assumption
and states the exact semantic resource used by the final candidate commit.
The separate \(\topts\)-safety condition similarly exposes a causal argument
that would be lost in a naive current-timestamp invariant.
\end{enumerate}

These strengthenings do not change the algorithm or the abstract
specification.  They make explicit the resources required for a stable,
compositional Hoare proof.

\appendix
\section{Correspondence with the mechanized definitions}

This appendix is only a navigation aid; none of the main argument depends on
Coq syntax.

\begin{center}
\begin{tabular}{>{\raggedright\arraybackslash}p{0.39\textwidth}
                >{\raggedright\arraybackslash}p{0.49\textwidth}}
\toprule
Mathematical notion & Mechanized name \\
\midrule
Concrete and branch invariants &
\code{concrete_wf}, \code{pool_represents}, \code{pool_protocol},
\code{timestamp_pending_edges}, \code{branch_represents} \\
Rectangular possibility family &
\code{pool_shared_eq}, \code{branch_merge},
\code{possibility_rectangular} \\
Global invariant & \code{I} \\
Rely and guarantee & \code{R}, \code{G} \\
Paper-style snapshot family & \code{GettingT} and the
\(\mapsto_{\exists},\mapsto_{\forall},\oplusset\) assertions \\
Candidate/row witnesses &
\code{candidate_view}, \code{row_snapshot_view},
\code{candidate_row_view} \\
Proof-relevant ignored-node elimination &
\code{node_cut_view}, \code{node_cuts_available} \\
Loop accumulator &
\code{EmptyEvidence}, \code{candidate_maximal},
\code{ScanAccumulator}, \code{ScanFold}, \code{ScanPending} \\
Candidate restriction and response &
\code{ac_actor_slice}, \code{getTop_candidate_commit_update} \\
Empty/failure protocol &
\code{CounterPending}, \code{getTop_counter_inv_update},
\code{getTop_counter_res_update} \\
Method triples &
\code{push_method_triple}, \code{getTop_method_triple},
\code{tryRemove_method_triple} \\
Layer theorem & \code{MListPool}, \code{MListPoolLinearizable} \\
\bottomrule
\end{tabular}
\end{center}

The source of the mechanized proof is
\path{examples/TSStack/ListPoolProof.v}; the implementation and specification
are in \path{examples/TSStack/ListPool.v} and
\path{examples/TSStack/ListPoolSpec.v}.

\end{document}
